\chapter{{Filtrations, processes and martingales}}\label{chap:filtration_martingale}


\section{Basic definitions}

First, recall the definitions of a filtration, an adapted process, a (sub)martingale, a stopping time and a stopped process, which are already in Mathlib.


\inputleannode{def:filtration}


\inputleannode{def:adapted}


\inputleannode{def:ProgMeasurable}


\inputleannode{lem:Adapted.progMeasurable_of_rightContinuous}




\inputleannode{def:Martingale}


\inputleannode{def:Submartingale}

\inputleannode{lem:condExp_sub_nonneg}




\inputleannode{lem:Submartingale.integrable_stoppedValue}




\inputleannode{lem:Martingale.congr}




\inputleannode{lem:Submartingale.congr}




\inputleannode{lem:conditional_jensen}




\inputleannode{cor:norm_condExp_le}




\inputleannode{lem:Martingale.submartingale_convex_comp}




\inputleannode{cor:Martingale.submartingale_norm}




\inputleannode{lem:convex_of_submg_is_submg}




\inputleannode{def:IsStoppingTime}


\inputleannode{def:StoppingTimeGen}


\inputleannode{lem:StoppingTimeGenMono}




\inputleannode{def:stoppedProcess}


\inputleannode{lem:Submartingale.stoppedProcess}




\inputleannode{def:hittingAfter}



\section{Right-continuous filtrations}

We now give the definition of a filtered probability space satisfying the usual conditions.


\begin{definition}[Left continuation]\label{def:leftLimitFiltration}
  \uses{def:filtration}
Assume that $T$ is a partial order. For $\mathcal{F}$ a filtration indexed by $T$ and $t \in T$, we define the \emph{left continuation} as
$$\mathcal{F}_{t-} =
  \begin{cases}
    \mathcal{F}_t, & \text{if $t$ is isolated on the left in the order topology}; \\
    \bigsqcup_{s < t} \mathcal{F}_s, & \text{otherwise}.
  \end{cases}
$$

Note that $\bigsqcup$ denotes the supremum in the lattice of sigma-algebras on $\Omega$.
\end{definition}


\inputleannode{def:rightLimitFiltration}


\inputleannode{def:rightLimitFiltration}




\inputleannode{lem:rightContDef}




\inputleannode{lem:rightContIsolated}




\inputleannode{lem:rightContSuccOrder}




\inputleannode{lem:rightContIsMax}




\inputleannode{lem:rightContExistsGt}




\inputleannode{lem:rightContNeBot}




\inputleannode{lem:rightContNotIsMax}




\inputleannode{lem:rightContEq}




\inputleannode{lem:leRightCont}




\inputleannode{lem:rightContSelf}




\begin{lemma}[Basic properties of the right continuation]\label{lem:rightLimitFiltration_basic}
  \uses{lem:filtrationRightCont, lem:rightContDef, lem:rightContIsolated, lem:rightContSuccOrder, lem:rightContIsMax, lem:rightContExistsGt,lem:rightContNeBot, lem:rightContNotIsMax, lem:rightContEq, lem:leRightCont, lem:rightContSelf}
  \leanok
Fake lemma for the dependency graph. Import this to depend on Definition~\ref{def:rightLimitFiltration}.
\end{lemma}




\inputleannode{def:rightContinuous}


\inputleannode{lem:eqRightCont}




\inputleannode{lem:rightContinuousRightCont}




\inputleannode{lem:rightContinuousMeasurableSet}




\begin{lemma}[Basic properties of right continuous filtrations]\label{lem:rightContinuous_basic}
  \uses{lem:eqRightCont, lem:rightContinuousRightCont, lem:rightContinuousMeasurableSet}
  \leanok
Fake lemma for the dependency graph. Import this to depend on Definition~\ref{def:rightContinuous}.
\end{lemma}




\inputleannode{def:hasUsualConditions}




\section{Predictable processes}


\inputleannode{def:predictableMeasurableSpace}


\inputleannode{def:predictable}


\inputleannode{lem:Predictable.progressive}




\inputleannode{lem:predictable_Ioc_prod}




\inputleannode{lem:predictable_nat_iff}




\section{Uniformly integrable}

\inputleannode{lem:condExpUI}




\inputleannode{lem:uniformIntegrable_stoppedValue_martingale}



\inputleannode{lem:uniformIntegrable_stoppedValue_martingale_of_countable_range}



\inputleannode{lem:uniformIntegrableAdd}



\inputleannode{lem:uniformIntegrableDominated}



\inputleannode{lem:uniformIntegrableDominatedSingleton}



\inputleannode{lem:uniformIntegrableNorm}



\inputleannode{lem:uniformIntegrableIffNorm}



\inputleannode{lem:uniformIntegrableComp}



\inputleannode{lem:uniformIntegrable_stoppedValue_submartingale}



\inputleannode{lem:memLp_of_tendstoInMeasure}



\inputleannode{lem:uniformIntegrable_of_tendstoInMeasure}



\inputleannode{lem:vitali}



\section{Optional sampling}

\inputleannode{def:isOrderedAddMonoid}

\inputleannode{def:isOrderedModule}

\inputleannode{def:orderClosedTopology}

\inputleannode{lem:optionalSampling_discrete}



\inputleannode{lem:optionalSampling_discrete_submartingale}



\inputleannode{lem:optionalSampling_discrete_supermartingale}



\inputleannode{def:approxSeq}

\inputleannode{def:approximableTimeIndex}

\inputleannode{lem:tendsto_stoppedValue_discreteApproxSequence}



\inputleannode{lem:discreteApproxSequence_of}



\inputleannode{lem:uniformIntegrable_stoppedValue_discreteApproxSequence}



\inputleannode{lem:tendsto_eLpNorm_stoppedValue_discreteApproxSequence}



\inputleannode{lem:stoppingTime_approximation}



\inputleannode{lem:stoppingTime_approximationNat}


\inputleannode{lem:optionalSampling}



\inputleannode{lem:optionalSamplingSubmartingale}



\section{Martingale convergence}

\inputleannode{def:limitProcess}


In Mathlib, we have results about convergence of martingales to their limit in discrete time.


\begin{theorem}\label{thm:tendsto_limitProcess_of_cadlag}
  \uses{def:Martingale, def:limitProcess}
Let $X$ be an uniformly integrable cadlag martingale with respect to the filtration $\mathcal{F}$.
Then there exists a limit process $X_\infty$ measurable with respect to $\mathcal{F}_\infty$ such that $X_t$ converges to $X_\infty$ almost surely as $t$ goes to infinity.
Furthermore, $X_t = P[X_\infty \mid \mathcal{F}_t]$ almost surely.
\end{theorem}







\section{Doob's Lp inequality}

In this section, we prove Doob's Lp inequality.


\inputleannode{lem:maximal_ineq}




\inputleannode{lem:doob_countable}




\begin{lemma}[Doob Lp Inequality for countable]\label{lem:doob_Lp_countable}
  Let $X : I \rightarrow \Omega \rightarrow \mathbb{R}$ be a non-negative sub-martingale. Let $I$ be countable.
  For every $M\in I,\lambda > 0$ and $p>1$ we have
  \begin{align*}
    \mathbb{E}\left[ \sup_{i\in I, i \leq M}X_i^p \right]
    \leq \left(\frac{p}{p-1}\right)^p\mathbb{E}[X_M^p]
    \: .
  \end{align*}
  That is, for $\Vert \cdot \Vert_p$ the $L^p$ norm,
  $\left\Vert \sup_{i \le M}  X_i  \right\Vert_p
    \leq \frac{p}{p-1} \left\Vert X_M \right\Vert_p
    \: .$
\end{lemma}

\begin{proof}
  \uses{lem:doob_countable}
\begin{align*}
  \mathbb{E}\left[ \sup_{i \le M}X_i^p \right]
  = p \int_0^\infty \mathbb{P}\left( \sup_{i \le M}X_i \geq \lambda \right) \lambda^{p-1} d\lambda
\end{align*}
By Theorem~\ref{lem:doob_countable} and then Fubini's theorem, we have then
\begin{align*}
  \mathbb{E}\left[ \sup_{i \le M}X_i^p \right]
  &\le p \int_0^\infty \mathbb{E}\left[X_M \mathbb{I}_{\sup_{i \le M}X_i \ge \lambda}\right] \lambda^{p-2} d\lambda
  \\
  &= p \mathbb{E}\left[X_M \int_0^{\sup_{i \le M}X_i} \lambda^{p-2} d\lambda\right]
  \\
  &= \frac{p}{p - 1} \mathbb{E}\left[X_M (\sup_{i \le M}X_i)^{p-1}\right]
  \: .
\end{align*}
Then by Hölder's inequality,
\begin{align*}
  \mathbb{E}\left[ \sup_{i \le M}X_i^p \right]
  &\le \frac{p}{p - 1} \left(\mathbb{E}[X_M^p]\right)^{1/p} \left(\mathbb{E}\left[\sup_{i \le M}X_i^p \right]\right)^{(p-1)/p}
  \: .
\end{align*}
We then divide the two sides by $\left(\mathbb{E}\left[\sup_{i \le M}X_i^p \right]\right)^{(p-1)/p}$ and raise to the power $p$ to conclude.
%  8.1.1 Pascucci.
\end{proof}


\inputleannode{thm:doob_ineq}




\inputleannode{cor:doob_ineq_norm}




\begin{theorem}[Doob's Lp inequality in $\mathbb{R}$]\label{thm:doob_lp}
  Let $X:\mathbb{R} \rightarrow \Omega \rightarrow \mathbb{R}$ be a right-continuous non-negative sub-martingale.
  For every $T, \lambda>0$ and $p>1$ we have
  \begin{align*}
    \mathbb{E}\left[ \sup_{t\in[0,T]}X_t^p \right]
    \leq \left(\frac{p}{p-1}\right)^p\mathbb{E}[X_T^p]
    \: .
  \end{align*}
  That is, for $\Vert \cdot \Vert_p$ the $L^p$ norm,
  $\left\Vert \sup_{t\in[0,T]}  X_t  \right\Vert_p
    \leq \frac{p}{p-1} \left\Vert X_T \right\Vert_p
    \: .$
\end{theorem}

\begin{proof}
  \uses{lem:doob_Lp_countable}
Since $X$ is right-continuous and $[0,T]$ is a compact interval, we have that
\begin{align*}
  \sup_{t\in[0,T]}X_t = \sup_{t\in[0,T] \cap \mathbb{Q}}X_t
  \: .
\end{align*}
Then apply Lemma~\ref{lem:doob_Lp_countable} with $I = [0,T] \cap \mathbb{Q}$ and $M = T$.
\end{proof}


\begin{corollary}[Doob's Lp inequality for normed spaces]\label{cor:doob_lp_norm}
  Let $X : \mathbb{R} \rightarrow \Omega\rightarrow E$ be a right-continuous martingale with values in a normed space $E$.
  For every $T, \lambda>0$ and $p>1$ we have
  \begin{align*}
    \mathbb{E}\left[ \sup_{t\in[0,T]} \lVert X_t \rVert ^ p \right]
    \leq \left(\frac{p}{p-1}\right)^p\mathbb{E}[\lVert X_T \rVert ^p]
    \: .
  \end{align*}
  That is, for $\Vert \cdot \Vert_p$ the $L^p$ norm,
  $\left\Vert \sup_{t\in[0,T]}  X_t  \right\Vert_p
    \leq \frac{p}{p-1} \left\Vert X_T \right\Vert_p
    \: .$
\end{corollary}

\begin{proof}
  \uses{cor:Martingale.submartingale_norm, thm:doob_lp}
  By Corollary~\ref{cor:Martingale.submartingale_norm}, $\lVert X \rVert$ is a sub-martingale.
  Then apply Theorem~\ref{thm:doob_lp}.
\end{proof}


\begin{lemma}[Stopped Submartingale]\label{lem:Submartingale.stoppedProcess_of_cadlag}
  \uses{def:Submartingale, def:stoppedProcess}
  Let $X:\mathbb{R}_+ \to \Omega \to \mathbb{R}$ be a cadlag submartingale and $\tau$ a stopping time.
  Then the stopped process $X^\tau$ is a submartingale.
\end{lemma}

\begin{proof}
  \uses{lem:Submartingale.stoppedProcess}

\end{proof}

\begin{lemma}[Doob Inequality for stopping times]\label{lem:doob_ineq_stop}
  Let $X:\mathbb{R}\times\Omega\rightarrow \mathbb{R}$ be a right-continuous non-negative sub-martingale.
  For every $\lambda>0$ and $p>1$ and $\tau$ stopping time a.s. bounded by $T>0$, we have
  $$
  P\left( \sup_{t\in[0,\tau]}X_t\geq\lambda \right)\leq \frac{\mathbb{E}[X_\tau]}{\lambda}.
  $$
\end{lemma}
\begin{proof}
  \uses{thm:doob_ineq, lem:Submartingale.stoppedProcess_of_cadlag}

\end{proof}

\begin{corollary}[Doob Inequality for stopping times in normed spaces]\label{cor:doob_ineq_stop}
  Let $X:\mathbb{R}\times\Omega\rightarrow E$ be a right-continuous martingale with values in a normed space $E$.
  For every $\lambda>0$ and $p>1$ and $\tau$ stopping time a.s. bounded by $T>0$, we have
  $$
  P\left( \sup_{t\in[0,\tau]}\lVert X_t \rVert \geq\lambda \right)\leq \frac{\mathbb{E}[\lVert X_\tau \rVert]}{\lambda}.
  $$
\end{corollary}
\begin{proof}
  \uses{cor:Martingale.submartingale_norm, lem:doob_ineq_stop}
  By Corollary~\ref{cor:Martingale.submartingale_norm}, $\lVert X \rVert$ is a sub-martingale.
  Then apply Theorem~\ref{lem:doob_ineq_stop}.
\end{proof}

\begin{lemma}[Doob's Lp Inequality for stopping times]\label{lem:doob_ineq_stop_exp_val}
  Let $X:\mathbb{R}\times\Omega\rightarrow \mathbb{R}$ be a right-continuous non-negative sub-martingale.
  For every $\lambda>0$ and $p>1$ and $\tau$ stopping time a.s. bounded by $T>0$, we have
  $$
  \mathbb{E}\left[ \sup_{t\in[0,\tau]}X_t^p \right]\leq \left(\frac{p}{p-1}\right)^p\mathbb{E}[X_\tau^p].
  $$
\end{lemma}
\begin{proof}
  \uses{lem:doob_ineq_stop, lem:Submartingale.stoppedProcess_of_cadlag}
  8.1.3 Pascucci.
\end{proof}

\begin{corollary}[Doob's Lp Inequality for stopping times in normed spaces]\label{cor:doob_ineq_stop_exp_val}
  Let $X:\mathbb{R}\times\Omega\rightarrow E$ be a right-continuous martingale with values in a normed space $E$.
  For every $\lambda>0$ and $p>1$ and $\tau$ stopping time a.s. bounded by $T>0$, we have
  $$
  \mathbb{E}\left[ \sup_{t\in[0,\tau]}\lVert X_t \rVert^p \right]\leq \left(\frac{p}{p-1}\right)^p\mathbb{E}[\lVert X_\tau \rVert^p].
  $$
\end{corollary}
\begin{proof}
  \uses{cor:Martingale.submartingale_norm, lem:doob_ineq_stop_exp_val}
  By Corollary~\ref{cor:Martingale.submartingale_norm}, $\lVert X \rVert$ is a sub-martingale.
  Then apply Theorem~\ref{lem:doob_ineq_stop_exp_val}.
\end{proof}
